%% LyX 2.4.4 created this file.  For more info, see https://www.lyx.org/.
%% Do not edit unless you really know what you are doing.
\documentclass[english]{article}
\usepackage[T1]{fontenc}
\usepackage[latin9]{inputenc}
\usepackage{enumitem}
\usepackage{amsmath}
\usepackage{amsthm}
\usepackage{cancel}
\usepackage{geometry}
\geometry{verbose,tmargin=1in,bmargin=1in,lmargin=1in,rmargin=1in}

\makeatletter
%%%%%%%%%%%%%%%%%%%%%%%%%%%%%% Textclass specific LaTeX commands.
\numberwithin{equation}{section}
\numberwithin{figure}{section}
\newlength{\lyxlabelwidth}      % auxiliary length 

%%%%%%%%%%%%%%%%%%%%%%%%%%%%%% User specified LaTeX commands.
\usepackage{fancyhdr}
\usepackage{lastpage}
\pagestyle{fancy}
\usepackage{enumitem}
\usepackage{ifthen}
%sets math fonts to be more readable
\everymath=\expandafter{\the\everymath\displaystyle}
%\DeclareMathSizes{10}{10}{10}{10}
\usepackage{multicol}

\usepackage{tikz}
\usetikzlibrary{plotmarks}
\usepackage{pgfplots}
\pgfplotsset{compat=newest}

\usepackage{calc}
\usepackage[nomessages]{fp}

\usepackage{advdate}    % Advancing/saving dates
\usepackage{datetime}   % Dates formatting
\usepackage{datenumber} % Counters for dates
% Added by lyx2lyx
\setlength{\parskip}{\smallskipamount}
\setlength{\parindent}{0pt}

\makeatother

\usepackage{babel}
\begin{document}
\renewcommand{\author}{Dr.\,James Ashe}

\newcommand{\classNum}{137}

\newcommand{\classSec}{B}

\newcommand{\nameType}{Name:}

\newcommand{\examType}{Final Exam}

\renewcommand{\date}{\formatdate{16}{5}{2013}}

%%First page's header%%%%%%%%%%%%%%%%%%%%%%
\fancypagestyle{firststyle} %%header settings for first pagr
{
	\fancyhead[L]{MTH \classNum{}(\classSec{}) \author{}\\
	\textbf{\examType{}} \date{}}
	\fancyhead[C]{\textbf{\nameType}}
	\setlength{\headsep}{14pt}
}
%%%%%%%%%%%%%%%%%%%%%%%%%%%%%%%%%%%%%%%%%%

%%Next page's header%%%%%%%%%%%%%%%%%%%%%%%
\setlist{nolistsep}
\lhead{\classNum{}(MTH \classSec{})} 
\chead{\textbf{\nameType}} 
\cfoot{\thepage{}~of~\pageref{LastPage}} 
\setlength{\headsep}{5pt} 
%%%%%%%%%%%%%%%%%%%%%%%%%%%%%%%%%%%%%%%%%%%

%%Various settings; see the preamble for other settings%%%%%
%%\setenumerate[0]{leftmargin=12pt,itemindent=0pt,labelwidth=0pt}
\renewcommand{\labelenumi}{\textbf{\arabic{enumi}.}}
\renewcommand{\labelenumii}{\textbf{\alph{enumii}.}}
\thispagestyle{firststyle}
%%%%%%%%%%%%%%%%%%%%%%%%%%%%%%%%%%%%%%%%%%
\newcommand{\xmax}{10}
\newcommand{\xmin}{-10}
\newcommand{\xstep}{0} %must be xmin+n
\newcommand{\ymax}{10}
\newcommand{\ymin}{-10}
\newcommand{\ystep}{0} %must be ymin+n

\newcommand{\dspace}{\vspace*{.75cm}}
\newcommand{\Dspace}{\dspace}

\textbf{You must show your work to receive credit. }Unless stated
otherwise, all problems are worth 5 points. 
\begin{enumerate}[resume]
\item {[}2pts each{]} Find the following:

\begin{enumerate}
\item $x^{5}-10000$ as $x\rightarrow\infty$\dspace
\item ${\displaystyle \lim_{x\rightarrow\infty}\frac{1}{x}}$\dspace
\item ${\displaystyle \lim_{x\rightarrow-3^{+}}\frac{2}{x+3}}$\dspace
\item The domain of $\frac{3x^{13}-12x+\sqrt{5}}{(x+3)(x-1)}$\dspace
\end{enumerate}
\end{enumerate}
\emph{For problems \ref{enu:solve-first} through \ref{enu:solve-last},
find all real solutions to the following equations. Write all answer
as exact values}\newcommand{\clfill}{\vspace{4cm}}
\begin{enumerate}[resume]
\item \begin{multicols*}{2}\label{enu:solve-first}$\frac{3x}{2}-1=2-\frac{7x}{4}$\vfill
\item $\frac{1}{2}|x-9|-7=9$\vfill
\item $x^{2}-11=0$\vfill\columnbreak
\item $\sqrt{x+3}=x+1$\vfill
\item $2e^{x-1}=10$\vfill
\item \label{enu:solve-last}$\log_{3}(4)=3-\log_{3}(x)$\vfill\columnbreak\end{multicols*}
\end{enumerate}
\emph{Solve the word problems. Approximate your answer to three decimal
places.}
\begin{enumerate}[resume]
\item Jeff knows that Holly paid \$40,230, including slaes tax for a new
Buick La-Sabre. If the sales tax is 10\%, then what is the cost of
the car \emph{before} the tax?\vfill
\item How long will it take \$100 to grow to \$350 if it is invested at
8\% compounded continuously?\vfill
\item Write an equation for the line shown in the graph. \\
\renewcommand{\xmax}{8} \renewcommand{\xstep}{4}
\renewcommand{\ymax}{\xmax} \renewcommand{\ystep}{4}
%%%%%%%
\begin{tikzpicture} 
	\begin{axis}[
		axis x line=middle, axis y line=middle,     		
		xtick={-\xmax,-\xstep,...,\xmax},
		ytick={-\ymax,-\ystep,...,\ymax},
		minor x tick num={3},
		minor y tick num={3},
		grid=both,
		xlabel=$x$,     
		ylabel={$x^2$},
		xmin=-\xmax.25,xmax=\xmax.25,
		ymin=-\ymax.25,ymax=\ymax.25,     		width=5cm,     		height=5cm] 		%\addplot[mark=noner,smooth,domain=-1:1]{}; 
		\addplot[color=red,thick, domain=-\xmax:\xmax] {-2/1*(x)-3};
	\end{axis}
\end{tikzpicture}
\item {[}2 pts each{]} Perform the indicated operations and write your answers
in the form $a+bi$.

\begin{enumerate}
\item $(2-i)(4+8i)$\vfill
\item $i^{4}$\dspace
\item $i^{25}$\dspace
\end{enumerate}
\item Let $f(x)=x+1$ and $g(x)=x^{2}-x$. Find and simplify the following
expressions.

\begin{enumerate}
\item {[}2 pts{]} $g(-2)$\dspace
\item $\left(g\circ f\right)(x)$\vfill\newpage
\end{enumerate}
\item {[}2pts each{]}\label{enu:quad-graph} Use the graph of the quadratic
function, to answer the following:

\begin{multicols}{2}
\begin{enumerate}
\item What is the vertex?\dspace
\item What is the range?\dspace
\item If any exist, find the $x$-int(s).\dspace
\item Find the $y$-int.\dspace
\item What are the zeros of $f(x)$?\dspace
\item Find the simplest polynomial with real coefficients that has the roots
from above (you do not have to expand).\columnbreak
\item []\renewcommand{\xmax}{10}
\renewcommand{\xmin}{-10}
\renewcommand{\xstep}{0} %must be xmin+n
\renewcommand{\ymax}{2}
\renewcommand{\ymin}{-18}
\renewcommand{\ystep}{0} %must be ymin+n
%%%%%%%
\begin{tikzpicture} 
\begin{axis}[height=6.5cm, 
	width=6.5cm, 
	axis lines=middle,
	minor x tick num={4},
	minor y tick num={4},
	grid=both,
	%xtick={0,50,...,250},
	%ytick={0,10,20,...,40},
	%axis line style={-latex}, 
	xmin=\xmin,xmax=\xmax, 
	ymin=\ymin,ymax=\ymax,
	%restrict y to domain=-1:10,
	%no markers, 
	xlabel={$x$},ylabel={$y$}, 
	%every axis x label/.append style={anchor=north}, 
	%every axis y label/.append style={anchor=east}
	]
	
	%%\addplot[color=red,solid,mark=*] coordinates {(4,-2)};
	\addplot[domain=\xmin:\xmax,thick,samples=100,draw=red,<->]{(x-2)^2-16} [yshift=0pt] node[pos=0.7,right] {$g(x)$};
	
\end{axis}
\end{tikzpicture}
\end{enumerate}
\end{multicols}
\item Consider the polynomials $P(x)=2x^{2}-3x+6$ and $(x-5)$.

\begin{enumerate}
\item Use division (synthetic or traditional) to find the quotient and remainder
when $P(x)$ is divided by $x-5$. \vfill\dspace
\item {[}2 pts{]} Using your answer above, find $P(5)$.\dspace
\item {[}2 pts{]} Is $x=5$ a root/zero of $P(x)$?\dspace
\end{enumerate}
\item {[}2 pts each{]} Use $f(x)=\frac{2x+1}{x+3}$

\begin{multicols}{2}
\begin{enumerate}
\item Find the domain of $f(x)$.\dspace
\item Find any vertical asymptotes.\dspace
\item Find the horizontal asymptote.\dspace\dspace
\item Sketch a graph of the function\columnbreak
\item []\renewcommand{\xmax}{10}
\renewcommand{\xmin}{-10}
\renewcommand{\xstep}{0} %must be xmin+n
\renewcommand{\ymax}{10}
\renewcommand{\ymin}{-10}
\renewcommand{\ystep}{0} %must be ymin+n
%%%%%%%
\begin{tikzpicture} 
\begin{axis}[height=5.5cm, 
	width=5.5cm, 
	axis lines=middle,
	minor x tick num={4},
	minor y tick num={4},
	grid=both,
	%xtick={0,50,...,250},
	%ytick={0,10,20,...,40},
	%axis line style={-latex}, 
	xmin=\xmin,xmax=\xmax, 
	ymin=\ymin,ymax=\ymax,
	%restrict y to domain=-1:10,
	no markers, 
	xlabel={$x$},ylabel={$y$}, 
	%every axis x label/.append style={anchor=north}, 
	%every axis y label/.append style={anchor=east}
	]
	
	%%\addplot[color=red,solid,mark=*] coordinates {(-3,-4)};
	%%\addplot[domain=\xmin:\xmax,thick,samples=100,draw=red,<->]{-(x-2)^2+5} node[right] {$$};
	
\end{axis}
\end{tikzpicture}
\end{enumerate}
\end{multicols}\newpage
\item Use $\frac{f(x+h)-f(x)}{h}$ for $h\neq0$, to complete and simplify
the following problem: Find the difference quotient of $f(x)=2x+1$.\vfill\dspace\dspace
\item {[}3 pts each{]} Evaluate the following:

\begin{enumerate}
\item $\log_{6}\left(6^{2}\right)$\dspace
\item $\frac{1}{3}e^{\ln(3)}$\dspace
\end{enumerate}
\item Let $g(x)=\log_{2}\left(x-2\right)$.

\begin{multicols}{2}
\begin{enumerate}
\item Graph $g(x)$.\dspace
\item {[}2 pts{]} Find the domain and range of $g(x)$.\dspace\columnbreak
\item []\renewcommand{\xmax}{10}
\renewcommand{\xmin}{-10}
\renewcommand{\xstep}{0} %must be xmin+n
\renewcommand{\ymax}{10}
\renewcommand{\ymin}{-10}
\renewcommand{\ystep}{0} %must be ymin+n
%%%%%%%
\begin{tikzpicture} 
\begin{axis}[height=5cm, 
	width=5cm, 
	axis lines=middle,
	minor x tick num={4},
	minor y tick num={4},
	grid=both,
	%xtick={0,50,...,250},
	%ytick={0,10,20,...,40},
	%axis line style={-latex}, 
	xmin=\xmin,xmax=\xmax, 
	ymin=\ymin,ymax=\ymax,
	%restrict y to domain=-1:10,
	no markers, 
	xlabel={$x$},ylabel={$y$}, 
	%every axis x label/.append style={anchor=north}, 
	%every axis y label/.append style={anchor=east}
	]
	
	%%\addplot[color=red,solid,mark=*] coordinates {(-3,-4)};
	%%\addplot[domain=\xmin:\xmax,thick,samples=100,draw=red,<->]{-(x+2)*(x-2)*(x-5)} node[pos=.65,right] {$f(x)$};
	
\end{axis}
\end{tikzpicture}
\end{enumerate}
\end{multicols}
\item Use the properties of logarithms to expand the expression as much
as possible.

\begin{enumerate}
\item []$\log_{11}\left(\frac{11}{y^{2}}\right)$\vfill
\end{enumerate}
\item Use the properties of logarithms to rewrite the expression as a single
logarithm.

\begin{enumerate}
\item []$\ln(4)+2\ln\left(3\right)$\vfill
\end{enumerate}
\item Find and sketch the graph of the equation of the parabola with focus
$(-1,6)$ and directrix $y=2$.\\
\renewcommand{\ymax}{14}
\renewcommand{\ymin}{-6}
\renewcommand{\xmax}{10}
\renewcommand{\xmin}{-10}
%%%%%%%
\begin{tikzpicture} 
\begin{axis}[height=5cm, 
	width=5cm, 
	axis lines=middle,
	minor x tick num={4},
	minor y tick num={4},
	%xtick={0,50,...,250},
	%ytick={0,10,20,...,40},
	%grid=both,
	%axis line style={-latex}, 
	xmin=\xmin,xmax=\xmax, 
	ymin=\ymin,ymax=\ymax,
	%restrict y to domain=-1:10,
	%no markers, 
	xlabel={$x$},ylabel={$y$}, 
	%every axis x label/.append style={anchor=north}, 
	%every axis y label/.append style={anchor=east}
	]
	
	%\addplot[color=black,solid,mark=*] coordinates {(-4,-4)};
	%\addplot[color=black,solid,mark=*] coordinates {(-4,-5)} node[below] {focus};
	%\addplot[domain=\xmin:\xmax,thick,samples=100,draw=blue,dashed]{-3} node[above,xshift=-1cm] {directrix:} node[below,xshift=-1cm,yshift=-.5cm] {$y=2$};
	%\addplot[domain=\xmin:\xmax,thick,samples=100,draw=red,<->]{-1/4*(x+4)^2-4} node[left] {$f(x)$};
	
\end{axis}
\end{tikzpicture}
\end{enumerate}

\end{document}
